\section{Method}\label{sec:method}
\paragraph{Pieces and paths.} A piece is a regular polygon $P$ with unit side and inradius $\rho$ (a unit square, triangle or hexagon), or a unit-diameter disk. Its \emph{shape state} is
\[
P(\gamma,\tau)=\gamma\big[(1-\tau)\,P\;\oplus\;\tau\rho\,D\big],
\]
the polygon scaled about its incenter, dilated by a disk of radius $\tau\rho$, and scaled by $\gamma$. At $\tau=1$ it is the inscribed disk, at $\tau=0$ the polygon, and for $\gamma\le1$ every state lies inside the polygon, so moving toward $(\gamma,\tau)=(1,0)$ only grows the piece. A \emph{path} is a schedule for $(\gamma,\tau)$:
\begin{itemize}\itemsep0pt
\item \textbf{harden}: $\gamma=1$, $\tau:1\to0$ (disks harden into polygons);
\item \textbf{grow}: $\tau=0$, $\gamma:\gamma_0\to1$ (rigid polygons grow, as in Lubachevsky--Stillinger \cite{lubachevsky});
\item \textbf{rigid}: $(\gamma,\tau)=(1,0)$ throughout, with the same schedule;
\end{itemize}
plus two classical baselines on rigid pieces: \textbf{sa}, Metropolis simulated annealing over poses and container size \cite{kirkpatrick}, and \textbf{pc}, perturbation--compression \cite{gensane} (shrink, relax, perturb on failure).

\paragraph{Energy.} The container (square, equilateral triangle or circle) is centred at the origin with size $L$. For poses $(x_i,y_i,\theta_i)$,
\[
E=\sum_{i<j}\big(2r-\operatorname{sd}_{ij}\big)_+^2+\sum_{i}\sum_{v}\big(\text{wall violation of }v\big)_+^2+\mu L,
\]
where $r=\gamma\tau\rho$, $\operatorname{sd}_{ij}$ is the signed distance between the cores (exact vertex--edge distance when apart, minus the separating-axis penetration depth when overlapping) and $v$ ranges over core vertices. The term $\mu L$ is inward pressure. Positions, angles and $L$ follow momentum gradient descent with analytic gradients (checked against finite differences, relative error below $10^{-7}$ on $1.7\times10^4$ components) and decaying Gaussian noise. A run compresses at pressure $3\mu$, follows its path while $L$ settles under pressure $\mu$, then lets $\mu$ decay with rigid pieces. In 3D the same holds for cubes with quaternion orientations \cite{nakajima12}.

\paragraph{From a run to a certified packing.} Penalties tolerate small overlaps, so every run is \emph{legalised}: centres are scaled apart about their mean until no pair overlaps by more than $10^{-12}$ under the separating-axis test, and the smallest container (translation free) is measured. Near-best runs are then \emph{tightened}: minimise $L$ over poses and one separating line (plane) per nearby pair, with every vertex inside the container, by SLSQP in trust-region outer iterations, keeping a step only if the independently legalised size decreases. A tightened packing below a catalogue value is spread to a clearance of $10^{-6}$ and checked twice: in floating point and in exact rational arithmetic (rational rotations from $\tan(\theta/2)$, pieces inflated by $10^{-8}$, an explicit rational separating direction for every pair).

\paragraph{Fair comparison.} All five methods are given the same number of pair evaluations per run (that of a \textbf{harden} run at the same $n$) and comparable tuning: $3\times3$ grids over their two main parameters on held-out instances, with fixed scoring rules (Section~\ref{sec:breadth}).
